\hojaejercicio{Hoja 1}

% Ejercicio 1
\ejercicio{enunciado}{no}{}{}

% Ejercicio 2
\ejercicio{enunciado}{no}{}{}

% Ejercicio 3
\ejercicio{enunciado}{no}{}{}

% Ejercicio 4
\ejercicio{enunciado}{no}{}{}

% Ejercicio 5
\ejercicio{enunciado}{no}{}{}

% Ejercicio 6
\ejercicio{enunciado}{no}{}{}

% Ejercicio 7
\ejercicio{enunciado}{no}{}{}

% Ejercicio 8
\ejercicio{enunciado}{no}{}{}

% Ejercicio 9
\ejercicio{enunciado}{no}{}{}

% Ejercicio 10
\ejercicio{enunciado}{no}{}{}

% Ejercicio 11
\ejercicio{enunciado}{no}{}{}

% Ejercicio 12
\ejercicio{enunciado}{no}{}{}

% Ejercicio 13
\ejercicio{enunciado}{no}{}{}

% Ejercicio 14
\ejercicio{enunciado}{no}{}{}

% Ejercicio 15
\ejercicio{enunciado}{no}{}{}

% Ejercicio 16
\ejercicio{enunciado}{no}{}{}

% Ejercicio 17
\ejercicio{resuelto}{profesor}{
    Sea $A \in M_n$ una matriz hermítica (resp. simétrica) definida positiva.
\begin{itemize}
    \item[a)] Probar que existe $B \in M_n$ hermítica (resp. simétrica) definida positiva tal que $A = B^2$.
    \item[b)] Demostrar que si $\text{cond}_2(A) > 1$ entonces $\text{cond}_2(B) < \text{cond}_2(A)$.
\end{itemize}
}
{
    \textbf{a)} Primero podemos aplicar el teorema de Schur para descomponer $A=UTU^*$, donde $U$ es una matriz unitaria (resp. ortogonal) y $T$ es triangular superior. Luego, como $A$ es h.d.p., $T$ debe ser diagonal con entradas positivas. Por lo tanto, podemos definir $\sqrt{D}$ como la matriz diagonal con las raíces cuadradas de las entradas de $D=T$:
    \[
        A=U\sqrt{D}\sqrt{D}U^*=(U\sqrt{D}U^*)^2 = B^2
    \]
    Por tanto, $B=U\sqrt{D}U^*$ es la matriz que buscamos, y es hermítica (resp. simétrica) y definida positiva, ya que $\sqrt{D}$ es diagonal con entradas positivas y $U$ es unitaria (resp. ortogonal).\\
    \textbf{b)} Ahora, como $U$ es unitaria (resp. ortogonal) y $\sqrt{D}$ es diagonal con entradas positivas, $B$ también es hermítica (resp. simétrica) y definida positiva, y además:
    \[
        \text{cond}_2(B) = 
        \| B\|_2 \| B^{-1} \|_2 = 
        \| U\sqrt{D}U^* \|_2 \| (U\sqrt{D}U^*)^{-1} \|_2 = 
        \| \sqrt{D} \|_2 \| \sqrt{D}^{-1} \|_2 = 
        \sqrt{\text{cond}_2(A)}
    \]
    Otra forma es darse cuenta de que:
    \[
        \| A \|_2^2 = \rho(A^*A) = \rho(A^2) = \rho^2(A) = \lambda_{\max}^2(A)
    \]
    Y la de su inversa:
    \[
        \| A^{-1} \|_2^2 = 
        \rho((A^{-1})^*A^{-1}) = 
        \rho((A^{-1})^2) = 
        \rho^2(A^{-1}) = 
        \lambda_{\max}^2(A^{-1}) =
        \lambda_{\min}^{-2}(A)
    \]
    Y por tanto su condición es: 
    \[
        \text{cond}_2(A) =
        \| A \|_2 \| A^{-1} \|_2 =
        \frac{\lambda_{\max}(A)}{\lambda_{\min}(A)}
    \]
    Análogamente para $B$, su condición es:
    \[
        \text{cond}_2(B) =
        \| B \|_2 \| B^{-1} \|_2 =
        \frac{\lambda_{\max}(B)}{\lambda_{\min}(B)} =
        \frac{\sqrt{\lambda_{\max}(A)}}{\sqrt{\lambda_{\min}(A)}} =
        \sqrt{\text{cond}_2(A)}
    \] 
    Por tanto, si $\text{cond}_2(A) > 1$ entonces $\text{cond}_2(B) < \text{cond}_2(A)$, ya que:
    \[
        \sqrt{\text{cond}_2(A)} < \text{cond}_2(A) \iff
        \text{cond}_2(A) > 1
    \]
}